\origpage{270}{ROBERTSON AND SEYMOUR}
\stmt
(7.9) \textit{Let $k \ge 0$ be an integer. Let $\mathscr{T}$ be a respectful tangle of order $\theta > k$ in a 2-cell drawing $\Gamma$ in a connected surface $\Sigma$; let its metric be $d$. Let $\Gamma'$ be a 2-cell subdrawing of $\Gamma$, and let $\mathscr{T}'$ be a respectful tangle of order $\theta - k$ in $\Gamma'$, with metric $d'$, such that $d'(a', b') \ge d(a, b) - k$ for all $a, b \in A(\Gamma)$ and $a', b' \in A(\Gamma')$ with $a \subseteq a'$ and $b \subseteq b'$. Then $(A \cap \Gamma', B \cap \Gamma') \in \mathscr{T}'$ for all $(A, B) \in \mathscr{T}$ of order $< \theta - k$. Moreover, if $\mathscr{T}''$ is another respectful tangle of order $\theta - k$ in $\Gamma'$, with metric $d''$, and again $d''(a', b') \ge d(a, b) - k$ for all $a, b \in A(\Gamma)$ and $a', b' \in A(\Gamma')$ with $a \subseteq a'$ and $b \subseteq b'$, then $\mathscr{T}'' = \mathscr{T}'$.}

\textit{Proof.} Let $\mathscr{T}_1$ be the tangle of order $\theta - k$ induced in $\Gamma$ by $\mathscr{T}'$, and let $\mathscr{T}_2$ be the set of members of $\mathscr{T}$ of order $< \theta - k$. To show that $(A \cap \Gamma', B \cap \Gamma') \in \mathscr{T}'$ for all $(A, B) \in \mathscr{T}$ of order $< \theta - k$, we must show that $\mathscr{T}_1 = \mathscr{T}_2$. Certainly $\mathscr{T}_1$ and $\mathscr{T}_2$ are both respectful tangles of order $\theta - k$ in $\Gamma$; let their metrics be $d_1$ and $d_2$, and let their inside functions be $\text{ins}_1, \text{ins}_2$. Let $K$ be a radial drawing of $\Gamma$, and suppose that $\mathscr{T}_1 \ne \mathscr{T}_2$. Then by [5, Theorem (6.3)], there is a circuit $C$ of $K$ of length $< 2(\theta - k)$ such that $\text{ins}_1(C) \ne \text{ins}_2(C)$. Choose $v \in V(C)$. By (1.2) there is a region $r$ of $K$ such that $d(\{v\}, r) = \theta$. Hence $d_2(\{v\}, r) = \theta - k$ by definition of $\mathscr{T}_2$, and so $r \notin \text{ins}_2(C)$; and $d_1(\{v\}, r) = \theta - k$, since $\{v\}$ and $r$ both correspond to atoms of $\Gamma'$, and so $r \notin \text{ins}_1(C)$. This is impossible, since $\text{ins}_1(C) \cup \text{ins}_2(C) = \Sigma$. We deduce that $\mathscr{T}_1 = \mathscr{T}_2$, and the first claim of the theorem holds.

For the second, let $\mathscr{T}''$ be as described in the theorem. Let $K'$ be a radial drawing of $\Gamma'$ and suppose that $\mathscr{T}' \ne \mathscr{T}''$. Then by [5, Theorem (6.3)] there is a circuit $C$ of $K'$ with length $< \theta - k$ such that $\text{ins}'(C) \ne \text{ins}''(C)$. Let $a$ be an atom of $\Gamma$ with $a \cap U(C) \ne \emptyset$. By (1.2) there is a region $r$ of $\Gamma$ with $d(a, r) = \theta$. Let $a', b' \in A(\Gamma')$ with $a \subseteq a'$ and $b \subseteq b'$; then
\begin{equation*}
d'(a', b'), d''(a', b') \ge d(a, r) - k = \theta - k
\end{equation*}
from the hypothesis. Since $C$ has length $< \theta - k$ it follows that $b' \notin \text{ins}'(C), \text{ins}''(C)$, a contradiction. Thus $\mathscr{T}' = \mathscr{T}''$. $\square$

We see in (7.8) that, while the order of $\mathscr{T}'$ is only $2k$ less than that of $\mathscr{T}$, the discrepancy in the distance functions $d(a, b) - d'(a', b')$ is larger, up to $4k + 2$. No doubt this is not best possible; perhaps the right answer is $2k$. But in view of (7.7) and (7.9), there are advantages in having these two discrepancies equal, and we therefore shall discard $2k + 2$ from the order of $\mathscr{T}'$, to obtain the following.

\stmt
(7.10) \textit{Let $k \ge 0$ be an integer, and let $\mathscr{T}$ be a respectful tangle of order $\ge 4k + 3$ in a 2-cell drawing $\Gamma$ in a connected surface $\Sigma$. Let its metric be $d$. Let $c \in A(\Gamma)$, and let $\Gamma'$ be a 2-cell subdrawing of $\Gamma$ such that $d(c, e) \le k$ for all $e \in E(\Gamma) - E(\Gamma')$, and such that if $c \in E(\Gamma) - E(\Gamma')$ then}
